% BANKS-SOURCE: pdf=76 print=66 kind=solution-section id=chapter06-implicit-solutions
\begin{bankssolutions}{Solutions to Implicit Exercises for Chapter 6}

% BANKS-SOURCE: pdf=76-77 print=66-67 kind=solution id=I-CH06-001
\BanksImplicitSolution{I-CH06-001}
Write the electron and muon traces as \(L^{\lambda\kappa}\) and
\(H_{\lambda\kappa}\).  Odd traces vanish, while the two- and four-gamma
traces give
\begin{align*}
 L^{\lambda\kappa}
 &=4\left[p_+^\lambda p_-^\kappa+p_+^\kappa p_-^\lambda
 -\eta^{\lambda\kappa}(p_+\!\cdot p_--m_e^2)\right],\\
 H_{\lambda\kappa}
 &=4\left[k_{-\lambda}k_{+\kappa}+k_{-\kappa}k_{+\lambda}
 -\eta_{\lambda\kappa}(k_-\!\cdot k_+-m_\mu^2)\right].
\end{align*}
The signs of both mass terms follow from the opposite signs of the particle
and antiparticle masses in each trace.  Contracting the tensors yields
\begin{align*}
 L^{\lambda\kappa}H_{\lambda\kappa}=32\bigl[&
 (p_-\!\cdot k_-)(p_+\!\cdot k_+)
 +(p_-\!\cdot k_+)(p_+\!\cdot k_-)
 -m_\mu^2(p_-\!\cdot p_+)\\
 &-m_e^2(k_-\!\cdot k_+)+2m_e^2m_\mu^2\bigr].
\end{align*}
The contraction is therefore \(32\) times its kinematic bracket.  With
\(\operatorname{tr}\mathbf1=4\), the passage from equation~\eqref{eq:6.7}
to the printed equation~\eqref{eq:6.8} loses this factor while retaining the
same prefactor.  Isolating the bracket requested here, its finite-electron-mass
addition is
\begin{equation*}
 \Delta\mathcal B_e
 =m_e^2\left(2m_\mu^2-k_-\!\cdot k_+\right)
 =\frac{m_e^2}{2}\left(2m_\mu^2-q^2\right),
\end{equation*}
where \(2k_-\!\cdot k_+=q^2+2m_\mu^2\).  This bracket-level result does not
adopt the inconsistent overall normalization of the printed equations.

In the center-of-mass frame \(q^2=-4E^2\), so
\(\Delta\mathcal B_e=m_e^2(2E^2+m_\mu^2)\).  The terms retained by Banks
are of order \(E^4\) at fixed angle.  Their ratio is therefore
\(O(m_e^2/E^2)\), which justifies the omission in the stated high-energy
regime.

% BANKS-SOURCE: pdf=77 print=67 kind=solution id=I-CH06-002
\BanksImplicitSolution{I-CH06-002}
\textit{(a) Weyl, Dirac, and Majorana bases.}
Use a second set of Pauli matrices \(\rho_i\) for the two-by-two block
index.  The Pauli identity
\(\sigma_i\sigma_j=\delta_{ij}\mathbf1+\ii\epsilon_{ijk}\sigma_k\)
gives the Weyl matrices
\begin{equation*}
 \gamma_W^0=\begin{pmatrix}0&\mathbf1\\ \mathbf1&0\end{pmatrix},
 \qquad
 \gamma_W^i=\begin{pmatrix}0&\sigma_i\\-\sigma_i&0\end{pmatrix},
 \qquad
 \gamma_{5,W}=\begin{pmatrix}-\mathbf1&0\\0&\mathbf1\end{pmatrix}.
\end{equation*}
Direct block multiplication proves
\(\{\gamma_W^\mu,\gamma_W^\nu\}=-2\eta^{\mu\nu}\).  It also shows that
\(\gamma_W^0\) is Hermitian and each \(\gamma_W^i\) is anti-Hermitian,
which is equivalent to
\begin{equation*}
 (\gamma_W^\mu)^\dagger=\gamma_W^0\gamma_W^\mu\gamma_W^0.
\end{equation*}

The unitary matrix
\begin{equation*}
 S_D=\frac{1}{\sqrt2}(\mathbf1-\ii\rho_2)\otimes\mathbf1
\end{equation*}
gives \(\gamma_D^\mu=S_D^\dagger\gamma_W^\mu S_D\), hence
\begin{equation*}
 \gamma_D^0=\rho_3\otimes\mathbf1
 =\begin{pmatrix}\mathbf1&0\\0&-\mathbf1\end{pmatrix},
 \qquad
 \gamma_D^i=\ii\rho_2\otimes\sigma_i
 =\begin{pmatrix}0&\sigma_i\\-\sigma_i&0\end{pmatrix}.
\end{equation*}
For the Majorana basis use the corrected unitary matrix adopted in the
transcription,
\begin{equation*}
 S_M=\frac{1}{\sqrt2}
 \begin{pmatrix}\mathbf1&\sigma_2\\ \sigma_2&-\mathbf1\end{pmatrix}
 =\frac{\rho_3\otimes\mathbf1+\rho_1\otimes\sigma_2}{\sqrt2}.
\end{equation*}
The two summands square to one and anticommute, so
\(S_M^\dagger S_M=1\).  Conjugating the Dirac matrices gives one useful
explicit Majorana set,
\begin{equation*}
 \gamma_M^0=\rho_1\otimes\sigma_2,
 \quad \gamma_M^1=\ii\mathbf1\otimes\sigma_3,
 \quad \gamma_M^2=-\ii\rho_2\otimes\sigma_2,
 \quad \gamma_M^3=-\ii\mathbf1\otimes\sigma_1.
\end{equation*}
Every entry is purely imaginary.  Consequently
\(-\ii\gamma_M^\mu\partial_\mu-m\) is a real differential operator, and
its equation holds separately for the real and imaginary parts of a
spinor.  Any unitary similarity transformation preserves both the Clifford
relations and the Hermiticity identity, since the adjacent factors of
\(SS^\dagger\) cancel.

\textit{(b) Antisymmetric basis and dualities.}
For the matrix basis, repeatedly anticommute gamma matrices until every word
is a sum of totally antisymmetric products.  Repeated indices then reduce by
the Clifford relation.  The numbers of products of ranks zero through four
are
\begin{equation*}
 \binom40+\binom41+\binom42+\binom43+\binom44
 =1+4+6+4+1=16.
\end{equation*}
Trace orthogonality between distinct antisymmetric grades proves linear
independence, so these sixteen matrices form a basis of all \(4\times4\)
matrices.

The phases in the duality relations require care.  Fix normalized
antisymmetrization by
\begin{equation*}
 \Gamma^{\mu_1\ldots\mu_r}
 :=\frac1{r!}\sum_{\pi\in S_r}\operatorname{sgn}(\pi)
 \gamma^{\mu_{\pi(1)}}\cdots\gamma^{\mu_{\pi(r)}}.
\end{equation*}
Evaluating the
\(0123\) component with \(\epsilon^{0123}=1\) fixes the coefficients:
\begin{align*}
 \Gamma^{\mu\nu\rho\sigma}
 &=-\ii\epsilon^{\mu\nu\rho\sigma}\gamma_5,\\
 \Gamma^{\mu\nu\rho}
 &=-\ii\epsilon^{\mu\nu\rho\sigma}\gamma_\sigma\gamma_5,\\
 \Gamma^{\mu\nu}
 &=\frac{\ii}{2}\epsilon^{\mu\nu\rho\sigma}
   \Gamma_{\rho\sigma}\gamma_5.
\end{align*}
The lower-rank formulas follow by contracting the rank-four identity with a
gamma matrix and using the Clifford relation.  Banks's Gordon convention is
\(\gamma^{\mu\nu}=-\Gamma^{\mu\nu}\), so the last identity retains the
same factor of \(\ii\).  Appendix~C prints these duality formulas without
the factors of \(\ii\).  Those phase-free forms cannot be combined with
its displayed definition
\(\gamma_5=\ii\gamma^0\gamma^1\gamma^2\gamma^3\): already the
\(0123\) component gives
\(\Gamma^{0123}=-\ii\gamma_5\).  The equations above are the identities
that follow from the stated definitions; a phase-free version requires a
corresponding redefinition of the antisymmetric products or of epsilon.

\textit{(c) Traces and contractions.}
Cyclicity and \(\{\gamma_5,\gamma^\mu\}=0\) imply that a trace of an odd
number of gamma matrices changes sign when \(\gamma_5^2\) is moved through
it.  Hence it vanishes.  The even traces begin with
\begin{align*}
 \operatorname{tr}\mathbf1&=4,
 &\operatorname{tr}(\gamma^\mu\gamma^\nu)&=-4\eta^{\mu\nu},\\
 \operatorname{tr}(\gamma^\mu\gamma^\nu\gamma^\rho\gamma^\sigma)
 &=4\left(\eta^{\mu\nu}\eta^{\rho\sigma}
 -\eta^{\mu\rho}\eta^{\nu\sigma}
 +\eta^{\mu\sigma}\eta^{\nu\rho}\right).
\end{align*}
The second formula follows by contracting its only possible rank-two tensor.
For the third, anticommute the first gamma successively to the right and use
cyclicity.  The same operation proves the recursion
\begin{equation*}
 \operatorname{tr}(\gamma^{\mu_1}\cdots\gamma^{\mu_{2r}})
 =\sum_{j=2}^{2r}(-1)^{j+1}\eta^{\mu_1\mu_j}
 \operatorname{tr}(\gamma^{\mu_2}\cdots
 \widehat{\gamma^{\mu_j}}\cdots\gamma^{\mu_{2r}}).
\end{equation*}
Here the hat means omission.  The basic \(\gamma_5\) traces are
\begin{equation*}
 \operatorname{tr}\gamma_5=0,
 \qquad
 \operatorname{tr}(\gamma_5\gamma^{\mu_1}\cdots\gamma^{\mu_r})=0
 \quad(1\leq r<4),
\end{equation*}
and
\begin{equation*}
 \operatorname{tr}(\gamma_5\gamma^\mu\gamma^\nu
 \gamma^\rho\gamma^\sigma)=-4\ii\epsilon^{\mu\nu\rho\sigma}.
\end{equation*}
Longer traces with \(\gamma_5\) reduce to this last expression by the same
pairing recursion.

Moving the outer gamma through a totally antisymmetric rank-\(r\) product
leaves \(4(-1)^{r+1}\Gamma_{\mu_1\ldots\mu_r}\) after the final gamma pair
contracts.  The \(r\) anticommutator terms sum, after antisymmetrization, to
\(2r(-1)^r\Gamma_{\mu_1\ldots\mu_r}\).  Adding them proves the compact
contraction identity
\begin{equation*}
 \gamma^\lambda\Gamma_{\mu_1\ldots\mu_r}\gamma_\lambda
 =(-1)^{r+1}(4-2r)\Gamma_{\mu_1\ldots\mu_r}.
\end{equation*}
For ordinary, unantisymmetrized strings this yields the Appendix~C examples
\begin{align*}
 \gamma^\lambda\gamma_\mu\gamma_\lambda&=2\gamma_\mu,\\
 \gamma^\lambda\gamma_\mu\gamma_\nu\gamma_\lambda&=4\eta_{\mu\nu},\\
 \gamma^\lambda\gamma_\mu\gamma_\nu\gamma_\kappa\gamma_\lambda
 &=-4\gamma_\mu\eta_{\nu\kappa}
   -2\gamma_\nu\gamma_\kappa\gamma_\mu,\\
 \gamma^\lambda\gamma_\mu\gamma_\nu\gamma_\kappa\gamma_\alpha
 \gamma_\lambda
 &=\gamma_\mu\left(4\gamma_\nu\eta_{\kappa\alpha}
   +2\gamma_\kappa\gamma_\alpha\gamma_\nu\right)
   -2\gamma_\nu\gamma_\kappa\gamma_\alpha\gamma_\mu.
\end{align*}
The third line uses the free-index correction in Erratum E-022.  In the last
line, the recurrence fixes the order
\(\gamma_\kappa\gamma_\alpha\gamma_\nu\).  Printed p.~250 instead reverses
the first two factors.  For four mutually distinct indices its right-hand
side vanishes, while direct Clifford reduction gives
\(-4\gamma_\mu\gamma_\nu\gamma_\kappa\gamma_\alpha\).
Expansion into antisymmetric grades extends these identities to every
length.

\textit{(d) Spinor solutions and normalization.}
For positive-frequency modes \(\ee^{+\ii p\cdot x}\), with
\(p\cdot x=-E_{\mathbf p}t+\mathbf p\cdot\mathbf x\), the kinetic operator
\(-\ii\slashed\partial-m\) gives
\((\slashed p-m)u=0\) and \((\slashed p+m)v=0\).
It remains to derive the spinors.  Put \(E_{\mathbf p}=\sqrt{\mathbf p^2+m^2}\),
\begin{equation*}
 -p_\mu\sigma^\mu=E_{\mathbf p}\mathbf1-\boldsymbol\sigma\!\cdot\mathbf p,
 \qquad
 -p_\mu\bar\sigma^\mu=E_{\mathbf p}\mathbf1+\boldsymbol\sigma\!\cdot\mathbf p.
\end{equation*}
These positive Hermitian matrices commute and their product is \(m^2\).
For orthonormal two-spinors \(\chi_s\) and \(\eta_s\), the Weyl-basis
solutions are
\begin{equation*}
 u_W(p,s)=\begin{pmatrix}+\sqrt{-p_\mu\sigma^\mu}\,\chi_s\\
 -\sqrt{-p_\mu\bar\sigma^\mu}\,\chi_s\end{pmatrix},
 \qquad
 v_W(p,s)=\begin{pmatrix}+\sqrt{-p_\mu\sigma^\mu}\,\eta_s\\
 +\sqrt{-p_\mu\bar\sigma^\mu}\,\eta_s\end{pmatrix}.
\end{equation*}
Substitution uses
\(\sqrt{-p_\mu\sigma^\mu}\sqrt{-p_\nu\bar\sigma^\nu}=m\), and proves
\((\slashed p-m)u=0\) and \((\slashed p+m)v=0\).  One convenient spin
convention is
\begin{equation*}
 \eta_s=2s\,\ii\sigma_2\chi_s^*,\qquad s=\pm\tfrac12.
\end{equation*}
The factor \(2s=\pm1\) fixes a spin-dependent phase and preserves
orthonormality.  It is present on source p.~250; the native appendix drops
the factor \(s\).

Applying \(S_D^\dagger\) to the displayed Weyl spinors, with a common
harmless phase chosen for \(u\), gives
\begin{equation*}
 u_D(p,s)=\begin{pmatrix}\dfrac{\boldsymbol\sigma\cdot\mathbf p}{\sqrt{E_{\mathbf p}+m}}\chi_s\\
 \sqrt{E_{\mathbf p}+m}\,\chi_s
 \end{pmatrix},
 \qquad
 v_D(p,s)=\begin{pmatrix}
 \sqrt{E_{\mathbf p}+m}\,\eta_s\\
 \dfrac{\boldsymbol\sigma\cdot\mathbf p}{\sqrt{E_{\mathbf p}+m}}\eta_s
 \end{pmatrix}.
\end{equation*}
Majorana-basis solutions are \(u_M=S_M^\dagger u_D\) and
\(v_M=S_M^\dagger v_D\).  Unitarity guarantees that their normalizations
are unchanged.  In every basis,
\begin{align*}
 u^\dagger(p,r)u(p,s)&=2E_{\mathbf p}\delta_{rs},
 &\bar u(p,r)u(p,s)&=-2m\delta_{rs},\\
 v^\dagger(p,r)v(p,s)&=2E_{\mathbf p}\delta_{rs},
 &\bar v(p,r)v(p,s)&=2m\delta_{rs}.
\end{align*}
Finally, using \(\sum_s\chi_s\chi_s^\dagger
=\sum_s\eta_s\eta_s^\dagger=\mathbf1\)
in the displayed block spinors gives
\begin{equation*}
 \sum_su(p,s)\bar u(p,s)=-(\slashed p+m),
 \qquad
 \sum_sv(p,s)\bar v(p,s)=-(\slashed p-m).
\end{equation*}
These are the completeness relations used in equation~\eqref{eq:6.7}.
They correspond to Banks's field-expansion factor
\((2E_{\mathbf p}(2\pi)^{D-1})^{-1/2}\); changing that factor would require a matching
rescaling of every spinor and spin sum.

% BANKS-SOURCE: pdf=78 print=68 kind=solution id=I-CH06-003
\BanksImplicitSolution{I-CH06-003}
In natural units an area has mass dimension \(-2\), so
\([\sigma_{\rm tot}]=-2\).  The QED coupling and
\(\alpha=e^2/(4\pi)\) are dimensionless in four dimensions.  The tree
amplitude contains two electromagnetic vertices, making the cross section
proportional to \(e^4\), or \(\alpha^2\).  When all masses vanish, \(E\)
is the only scale.  Dimensional analysis then forces
\begin{equation*}
 \sigma_{\rm tot}=C\frac{\alpha^2}{E^2}.
\end{equation*}

Taking equation~\eqref{eq:6.11} as the normalization stated in the exercise,
set \(a=m_\mu^2/E^2\) and \(b=\mu^2/(4E^2)\).  It becomes
\begin{equation*}
 \sigma_{\rm tot}=\frac{\pi\alpha^2}{3E^2}
 (1-b)^{-2}(1-a)^{1/2}\left(1+\frac a2\right).
\end{equation*}
Thus \(C=\pi/3\).  Expanding the two dimensionless factors gives
\begin{align*}
 (1-b)^{-2}&=1+2b+3b^2+O(b^3),\\
 (1-a)^{1/2}\left(1+\frac a2\right)
 &=1-\frac{3a^2}{8}+O(a^3).
\end{align*}
Consequently
\begin{equation*}
 \sigma_{\rm tot}=\frac{\pi\alpha^2}{3E^2}
 \left[1+\frac{\mu^2}{2E^2}
 -\frac{3m_\mu^4}{8E^4}
 +O\!\left(\frac{\mu^4}{E^4},
 \frac{\mu^2m_\mu^4}{E^6},\frac{m_\mu^6}{E^6}\right)\right].
\end{equation*}
The term linear in \(m_\mu^2/E^2\) cancels between phase space and the
matrix element.  A nonzero electron mass adds dimensionless corrections
starting at \(m_e^2/E^2\), as the preceding implicit exercise shows.  At
higher loop order, logarithms of mass ratios can multiply these allowed
powers without changing the overall dimension \(E^{-2}\).

% BANKS-SOURCE: pdf=79 print=69 kind=solution id=I-CH06-004
\BanksImplicitSolution{I-CH06-004}
Let a dot now mean \(d/dv\).  The substitutions
\begin{equation*}
 s=\frac{u}{M_\mu^2},\qquad
 \tau=\frac{v}{M_\mu^2},\qquad
 d\tau=\frac{dv}{M_\mu^2},\qquad
 \frac{dx^\mu}{d\tau}=M_\mu^2\dot x^\mu
\end{equation*}
separate the four contributions to the exponent of
equation~\eqref{eq:6.15}:
\begin{align*}
 -\ii sM_\mu^2&=-\ii u,\\
 \ii\int_0^s d\tau\left(\frac{dx}{2d\tau}\right)^2
 &=\frac{\ii M_\mu^2}{4}\int_0^u dv\,\dot x^2,\\
 \ii\int_0^s d\tau\,A_\mu(x)\frac{dx^\mu}{d\tau}
 &=\ii\int_0^u dv\,A_\mu(x)\dot x^\mu,\\
 \frac{\ii}{2}\int_0^s d\tau\,\sigma^{\mu\nu}F_{\mu\nu}(x)
 &=\frac{\ii}{2M_\mu^2}\int_0^u dv\,
   \sigma^{\mu\nu}F_{\mu\nu}(x).
\end{align*}
For fixed endpoints, \(u=O(1)\), and backgrounds held fixed as the mass
grows, these terms are respectively \(O(1)\), \(O(M_\mu^2)\),
\(O(1)\), and \(O(M_\mu^{-2})\).  The dimensions provide a quick check:
\([A]=1\), \([F]=2\), and a line integral of \(A\) is dimensionless.

Varying the leading kinetic functional gives \(\ddot x^\mu=0\).  With
\(x(0)=x\) and \(x(u)=y\), its saddle is
\begin{equation*}
 x_{\rm cl}(v)=x+\frac vu(y-x),
\end{equation*}
the straight path used immediately after equation~\eqref{eq:6.15}.  The
vector coupling is smaller than the kinetic action by \(M_\mu^{-2}\), so
it displaces this saddle at that relative order while contributing an
order-one Wilson-line phase.  The spin interaction is itself
\(O(M_\mu^{-2})\); compared directly with the kinetic action it is down by
\(M_\mu^{-4}\).  This counting assumes that the background amplitude and
its variation scale stay independent of \(M_\mu\), precisely the slowly
varying heavy-particle regime described by Banks.

% BANKS-SOURCE: pdf=81-82 print=71-72 kind=solution id=I-CH06-005
\BanksImplicitSolution{I-CH06-005}
Put \(P=p'+p=2p+q\).  After all contracted factors of \(\slashed p\) and
\(\slashed p'\) have been moved next to the external spinors, a general
parity-even vector can be written between those spinors as a linear
combination of
\begin{equation*}
 \gamma_\mu,\quad P_\mu\mathbf1,\quad q_\mu\mathbf1,
 \quad\gamma_{\mu\nu}P^\nu,\quad\gamma_{\mu\nu}q^\nu.
\end{equation*}
Terms containing \(\gamma_5\), \(\gamma_\mu\gamma_5\), or an epsilon
tensor are parity odd for this vector current.  Extra slashed momenta add no
covariants because the on-shell equations reduce them after anticommutation.
For example,
\begin{equation*}
 \bar u(p')\gamma_\mu\slashed q\,u(p)
 =-2p'_\mu\bar u(p')u(p)-2m\bar u(p')\gamma_\mu u(p),
\end{equation*}
and reversing the order of the two matrices gives the analogous combination.

The equal-mass Dirac equations first give
\begin{equation*}
 \bar u(p')\slashed q\,u(p)
 =\bar u(p')(\slashed p'-\slashed p)u(p)=0.
\end{equation*}
With Banks's convention
\(\gamma^{\mu\nu}=\tfrac12(\gamma^\nu\gamma^\mu-
\gamma^\mu\gamma^\nu)\), the Gordon identity reads
\begin{equation*}
 P_\mu\bar u(p')u(p)
 =\bar u(p')\left(\gamma_{\mu\nu}q^\nu-2m\gamma_\mu\right)u(p).
\end{equation*}
Equivalently,
\(2m\bar u\gamma_\mu u=\bar u(\gamma_{\mu\nu}q^\nu-P_\mu)u\),
which is the convention used in Problem~5.14.  This form fixes the tensor
term and eliminates \(P_\mu\mathbf1\).  A second direct use of the Dirac
equations gives
\begin{align*}
 \bar u(p')\gamma_{\mu\nu}P^\nu u(p)
 &=\frac12\bar u(p')\left(\slashed P\gamma_\mu
 -\gamma_\mu\slashed P\right)u(p)\\
 &=+q_\mu\bar u(p')u(p),
\end{align*}
so \(\gamma_{\mu\nu}P^\nu\) is also redundant.  These identities leave only
\(\gamma_\mu\), \(\gamma_{\mu\nu}q^\nu\), and
\(q_\mu\mathbf1\).  Absorbing convenient powers of \(2m\) into scalar
coefficients gives
\begin{equation*}
 \bar u(p')\Gamma_\mu u(p)=\bar u(p')\left[
 F_1\gamma_\mu+F_2\frac{\gamma_{\mu\nu}q^\nu}{2m}
 +F_3q_\mu\right]u(p).
\end{equation*}
On shell, \(p^2=p'^2=-m^2\) and
\(2p\!\cdot q=-q^2\).  Every scalar coefficient can therefore depend
only on \(q^2\) and \(m^2\).  The three displayed covariants are independent
before current conservation is imposed.

% BANKS-SOURCE: pdf=82-83 print=72-73 kind=solution id=I-CH06-006
\BanksImplicitSolution{I-CH06-006}
Contract the proposed basis with \(q^\mu\).  On shell,
\begin{equation*}
 q\!\cdot(2p+q)=(p+q)^2-p^2=0,
 \qquad
 \bar u(p+q)\slashed q\,u(p)=0.
\end{equation*}
Current conservation therefore reduces to
\begin{equation*}
 0=q^\mu\bar u\Gamma_\mu u
 =C(q^2)q^2\bar u(p+q)u(p).
\end{equation*}
For nonzero \(q^2\), spin states can be chosen with a nonzero scalar
bilinear, so \(C(q^2)=0\).  The regulated one-particle-irreducible vertex
has no \(1/q^2\) pole at the origin.  Continuity from a punctured
neighborhood then fixes \(C(0)=0\) as well.  Equivalently, impose the Ward
identity before taking the constant-field limit \(q\to0\).  This explains
why the identity at the single point \(q^2=0\) must not be read in
isolation.

The explicit parameter integral supplies an independent check.  Its
longitudinal numerator is proportional to
\begin{equation*}
 q_\mu m(z-2)(x-y).
\end{equation*}
The simplex measure
\(dx\,dy\,dz\,\delta(1-x-y-z)\) is invariant under
\(x\leftrightarrow y\), and the shifted denominator depends on these
variables through \(xy\) and \(x+y=1-z\).  The rest of the integrand is
symmetric, whereas \(x-y\) is odd.  Its integral vanishes for every
\(q^2\), including zero, in agreement with current conservation.

% BANKS-SOURCE: pdf=84 print=74 kind=solution id=I-CH06-007
\BanksImplicitSolution{I-CH06-007}
The dimensionless integral multiplying the magnetic covariant in
equation~\eqref{eq:6.18} is
\begin{equation*}
 I(q^2)=\int_0^1dx\,dy\,dz\,\delta(1-x-y-z)
 \frac{m^2z(1-z)}{(1-z)^2m^2+xyq^2}.
\end{equation*}
At \(q^2=0\), integrate over \(y\) and use
\(0\leq x\leq1-z\):
\begin{align*}
 I(0)
 &=\int_0^1dz\int_0^{1-z}dx\,
   \frac{z}{1-z}\\
 &=\int_0^1dz\,z=\frac12.
\end{align*}
The factor \(1-z\) from the length of the \(x\)-interval cancels the
apparent endpoint singularity.  Thus the endpoint at \(z=1\) is integrable.
Keeping a positive photon-mass regulator through the Euclidean integral
replaces the denominator at \(q^2=0\) by
\(m^2(1-z)^2+\mu^2z\).  This plus sign follows by completing the square in
the three Feynman denominators; the nearby printed formula for \(U\) has the
opposite sign on \(z\mu^2\).  After the \(x\) integration, its integrand is
\begin{equation*}
 f_\mu(z)=\frac{m^2z(1-z)^2}{m^2(1-z)^2+\mu^2z},
 \qquad 0\leq f_\mu(z)\leq z.
\end{equation*}
Dominated convergence gives
\(\lim_{\mu\to0}\int_0^1 f_\mu(z)\,dz=1/2\), which justifies taking the
massless-photon limit.

Equation~\eqref{eq:6.18} already multiplies the same covariant
\(+\gamma^{\mu\nu}q_\nu/(2m)\) used to define \(F_2\).  Its coefficient is
therefore
\begin{equation*}
 F_2(0)=\frac{e^2}{4\pi^2}I(0)
 =\frac{e^2}{8\pi^2}=\frac{\alpha}{2\pi}.
\end{equation*}
The magnetic covariant starts at first order in the external momentum, which
improves the soft behavior of the loop integral.  The electric form factor
lacks that improvement and retains the infrared divergence discussed by
Banks.

\end{bankssolutions}
